% Generated draft: requires mathematical and source review.
\documentclass[11pt]{article}
\usepackage[margin=25mm]{geometry}
\usepackage{fontspec}
\setmainfont{Latin Modern Roman}
\usepackage{amsmath,amssymb,mathtools,unicode-math}
\setmathfont{Latin Modern Math}
\usepackage{enumitem,needspace,xurl,hyperref}
\setlength{\parindent}{0pt}
\setlength{\parskip}{5pt}
\setlength{\emergencystretch}{3em}
\newcommand{\N}{\mathbb{N}}
\newcommand{\Z}{\mathbb{Z}}
\newcommand{\Q}{\mathbb{Q}}
\newcommand{\R}{\mathbb{R}}
\newcommand{\C}{\mathbb{C}}
\newcommand{\F}{\mathbb{F}}
\newcommand{\A}{\mathbb{A}}
\newcommand{\PP}{\mathbb{P}}
\newcommand{\E}{\mathbb{E}}
\newcommand{\eps}{\varepsilon}
\newcommand{\dd}{\,\mathrm{d}}
\begin{document}
\section*{UnsolvedMath ID: 30003718}
\textbf{Source problem code:} catalogue entry
\subsection{Definitions and mathematical statement}
Let \((M,g)\) be a Riemannian manifold, with Riemannian distance \(d_g\). The supplied record does not specify the data model, geometric hypotheses on \(M\), or a definition of a principal flow. Consequently, the following is a formulation of the supplied research questions with their unresolved specifications made explicit, rather than a fully quantified conjecture.

For terminology only, if the data are represented by a Borel probability measure \(\mu\) on \(M\), the usual squared-distance Fr\'echet potential is
\[
F_\mu(p)=\int_M d_g(p,x)^2\,\mathrm{d}\mu(x),
\]
where this integral is finite. A local Fr\'echet mean is a point \(p\) for which there exists a neighborhood \(U\) of \(p\) such that \(F_\mu(p)\leq F_\mu(q)\) for every \(q\in U\). This conventional definition supplies terminology; the record does not specify whether the intended potential is this function, an empirical version, or a scale-dependent localized potential.

Write \(h\) for a proposed scale parameter and \(\gamma_h:I_h\to M\) for a proposed principal-flow curve, with its defining rule and parameter domain still to be specified. A parametrized curve has a self-intersection if \(\gamma_h(s)=\gamma_h(t)\) for some distinct \(s,t\in I_h\); whether repeated traversal or closure is to count as an unwanted intersection is not specified. For a specified real-valued potential \(F\), its sublevel sets are \(\{p\in M:F(p)\leq a\}\), indexed by \(a\in\R\). Persistence diagrams associated with this filtration require choices of homology degree, coefficients, and hypotheses ensuring the intended diagram is defined. No such choices are supplied.

\begin{enumerate}
\item Determine whether a scale parameter can describe a principal flow that ``dies out.'' The meaning of dying out, the variable with respect to which it occurs, and the relevant endpoint or limiting regime must be specified.
\item Determine whether additional regularization is needed to prevent self-intersection or winding of principal flows. The underlying construction, admissible regularizers, and meaning of unwanted winding must be specified before a necessity or sufficiency assertion can be stated.
\item Determine whether persistence diagrams of the intended Fr\'echet potential can detect local means at different scales. The relationship between the persistence filtration parameter and the geometric scale, and the meaning of detection, remain unspecified.
\item Determine whether one scale parameter can govern both the choice of starting points and the estimation of local covariance. The starting-point rule, localization procedure, and tangent-space covariance construction remain unspecified.
\item Develop higher-dimensional analogues of principal flows, including canonical principal submanifolds and related generalizations. The dimension, admissible class of submanifolds, defining criterion, and meaning of canonical must be supplied.
\end{enumerate}

The source gives no universal or existential quantifiers over manifolds, data sets or measures, scales, or regularization classes. It also supplies no boundary conditions, completeness assumptions, stopping rules, or uniqueness requirements. These omissions prevent replacing the research questions by a single precise theorem or conjecture without adding hypotheses.
\subsection{Short English statement}
The problem asks for a multiscale theory of principal flows on Riemannian manifolds: how flows terminate, how to control intersections and winding, how persistence might identify local Fr\'echet means, whether one scale can govern initialization and local covariance, and how to obtain canonical higher-dimensional principal submanifolds. The source leaves the constructions and success criteria unspecified, so this draft preserves the questions and identifies the definitions needed for a precise conjecture.
\subsection{Sources}
\begin{itemize}
\item Local UnsolvedMath registry, supplied complete record: permanent ID \(30003718\), title \emph{Multiscale Principal Flows on Riemannian Manifolds}, exact source problem code \texttt{catalogue entry}; collection: \emph{Oberwolfach Reports Open Problems}. No individual proposer is supplied; the proposed year is recorded as 2018.
\item Supplied report citation: \emph{Statistics for Data with Geometric Structure (2019)}, Oberwolfach Reports, pp.\ 125--186, DOI \texttt{10.4171/owr/2018/3}, \url{https://doi.org/10.4171/owr/2018/3}. The background attributes the questions to mathematical workshop discussions documented in Oberwolfach Reports. The supplied title year is 2019, whereas the DOI and report label refer to 2018; this discrepancy is unresolved here.
\item Supplied literature source: \emph{Oberwolfach Report 3/2018}, \url{https://doi.org/10.4171/owr/2018/3}. The literature assessment was checked on \texttt{2026-08-21}, according to the supplied record. It labels the problem open and describes the questions about dying-out scales, self-intersection, multiscale local means, and higher-dimensional principal flows as broad unresolved problems without a general canonical construction. This is supplied dated curation derived from the OWR material, not an independently verified assessment; time-sensitive status and cited claims require independent verification.
\item The registry labels the statement \texttt{corrected\_verified} and says that notation, typography, and sentence structure were repaired while preserving mathematical content. These are supplied editorial labels and claims, not mathematical verification of this draft or independent confirmation of the cited report.
\end{itemize}
\end{document}
